\section{Basic finite energy potentials on a Hopf manifold}
\label{sec:hopf-basic-perron}

The Hopf geometry of \cref{sec:hopf-all-dimensions} gives a nonclosed background for which a distinguished family of potentials has a complete geodesic metric. Its Perron envelope uses all subgeodesics, although the metric measures lengths within the distinguished family. The distinction matters: allowing additional smooth paths can decrease the length.

Fix $n\geq3$ and put $d=n-1$. Write
\[
 X=H_n,\qquad Y=\mathbb P^d,\qquad \pi=p_n,
 \qquad \omega=\omega_n=\alpha_n+\beta_n.
\]
Let $\kappa$ be the mass-one Fubini--Study form on $Y$, so $\alpha_n=\pi^*\kappa$. On the cover, $r=\log|z|^2$ is defined modulo two and $\beta_n=\mathrm dr\wedge\dc r$. The compact elliptic group
\[
 E=\mathbb C/(\mathbb Z+2\pi\mathrm i\mathbb Z)
\]
acts by scalar multiplication, preserves $\omega$, and acts transitively on each fiber. A potential is \emph{basic} if it is invariant under this action; such a potential is exactly $\pi^*v$.\index{Hopf manifold!basic potentials}

These manifolds are not products of positive-dimensional compact complex manifolds. Indeed, $H_n$ is diffeomorphic to $S^{2n-1}\times S^1$, whose intermediate positive even-degree real cohomology vanishes. A product with a compact complex factor of dimension $k$, $0<k<n$, would have a nonzero degree-$2k$ class: pullback of that factor's top class is injective because projection has a smooth section. Thus the fiber reduction below takes place on a nonproduct manifold.

\Needspace{14\baselineskip}
\begin{proposition}\label{prop:hopf-basic-descent}
Let $S$ be a strip and $V$ a genuine quasi-plurisubharmonic function on $Y\times S$. Then
\[
 \operatorname{pr}_X^*\omega+\ddc(\pi\times\mathrm{id})^*V\geq0
 \quad\Longleftrightarrow\quad
 \operatorname{pr}_Y^*\kappa+\ddc V\geq0.
\]
For every $v\in\PSH(Y,\kappa)$, the extended non-pluripolar measure satisfies
\begin{equation}\label{eq:hopf-basic-measure}
 \MA^{\mathrm{ext}}_\omega(\pi^*v)
   =n\beta_n\wedge\pi^*\MA_\kappa(v).
\end{equation}
In particular, for every nonnegative Borel function $F$,
\begin{equation}\label{eq:hopf-basic-fiber-integral}
 \int_X\pi^*F\,\MA^{\mathrm{ext}}_\omega(\pi^*v)
 =2n\int_YF\,\MA_\kappa(v).
\end{equation}
If the base Perron family between $v_0,v_1$ is nonempty, the full Hermitian Perron envelope between $\pi^*v_0,\pi^*v_1$ is the pullback of the base envelope.
\end{proposition}
\begin{proof}
Write $\alpha=\alpha_n$ and $\beta=\beta_n$. Locally
$\beta=(\mathrm i/(2\pi))\partial r\wedge\bar\partial r$, and $\partial r$ has a nonzero vertical component. Every vector on $Y\times S$ therefore has a lift annihilated by $\partial r$. Testing the current on these lifts removes $\beta$; integrating the resulting distributional inequality against a nonnegative compact fiber density gives positivity on the base. Conversely, pullback of a positive base current plus $\beta$ is positive. An invariant genuine potential descends canonically by restricting to local holomorphic sections, so the argument also applies to singular potentials.

For bounded $v$, the local product with smooth corrections from \cref{chap:products} gives
\[
 (\omega+\ddc\pi^*v)^n
  =(\beta+\pi^*(\kappa+\ddc v))^n
  =n\beta\wedge\pi^*(\kappa+\ddc v)^d.
\]
All other terms vanish by $\beta^2=0$ and the dimension of $Y$. Apply this identity to $\max(v,-j)$ and restrict to $\{v>-j\}$. Increasing finite-level limits give \eqref{eq:hopf-basic-measure}, without assigning mass to the pole set. The identity $\int_{\pi^{-1}(y)}\beta=2$ from \cref{sec:hopf-all-dimensions} gives \eqref{eq:hopf-basic-fiber-integral}.

The full Perron competitor family is preserved by $E$. Its supremum and upper semicontinuous regularization are invariant. The chord bound and a pulled-back base competitor make this envelope a genuine subgeodesic by \cref{prop:hermitian-perron}. Positivity descent makes its descended envelope a base competitor with the prescribed joint boundary upper bounds. Conversely, every base competitor pulls back to a full competitor. The two envelopes consequently agree; in submersion charts, upper semicontinuous regularization commutes with pullback.
\end{proof}

Let $\mathcal E^2(Y,\kappa)$ denote the canonical $\kappa$-plurisubharmonic potentials of non-pluripolar mass one and finite squared energy. Define
\[
 \mathcal B_\infty=\{\pi^*v:v\in C^\infty(Y),\ \kappa+\ddc v>0\},
 \qquad
 \mathcal B_2=\pi^*\mathcal E^2(Y,\kappa).
\]
Thus $\mathcal B_2$ consists precisely of the basic potentials of extended mass $2n$ with $\int_X|u|^2\MA^{\mathrm{ext}}_\omega(u)<\infty$. On $\mathcal B_\infty$, let $d_B$ be the infimum of the lengths of smooth basic paths for the metric $G$ above.\index{Finite energy!basic Hopf potentials}

We use the classical K\"ahler completion theorem: the $L^2$ length metric for $\int_Y\xi\eta(\kappa+\ddc v)^d$ has completion $\mathcal E^2(Y,\kappa)$, is a CAT$(0)$ metric, and its geodesic between two points is their Perron segment. That segment has capacity endpoint limits and all its slices lie in $\mathcal E^2$. We also use the full-mass characterization $P_\kappa[v]=0$ for $v\in\mathcal E^2$. These are closed-background inputs on $Y$; their exact source and operator normalization are specified in \cref{history:hopf-basic-perron}.

\Needspace{12\baselineskip}
\begin{theorem}\label{thm:hopf-basic-completion}\index{Metric completion!basic Hopf potentials}
The completion of $(\mathcal B_\infty,d_B)$ is $\mathcal B_2$. It is a CAT$(0)$ space, and its unique constant-speed geodesic between two points is their full Hermitian Perron segment. This segment attains both endpoints in $L^1(X,\mathrm dV)$ for every smooth positive volume form $\mathrm dV$. For every $u\in\mathcal B_2$ and $c\in\mathbb R$,
\[
 d_B(u,u+c)=\sqrt{2n}|c|.
\]
\end{theorem}
\begin{proof}
Fiber integration gives the exact restricted metric
\begin{equation}\label{eq:hopf-basic-metric}
 G_{\pi^*v}(\pi^*\xi,\pi^*\eta)
 =2n\int_Y\xi\eta(\kappa+\ddc v)^d.
\end{equation}
The K\"ahler completion theorem therefore gives the stated completion and CAT$(0)$ property after multiplying base distances by $\sqrt{2n}$. If $\widetilde d_D$ denotes the completed metric written with operator $\mathrm i\partial\bar\partial$ and background $\kappa$, the normalization is explicitly
\begin{equation}\label{eq:hopf-basic-distance}
 d_B(\pi^*v,\pi^*w)
 =2\pi\sqrt{2n}\,\widetilde d_D(v/(2\pi),w/(2\pi)),
\end{equation}
because $\kappa+\mathrm i\partial\bar\partial(v/(2\pi))=\kappa+\ddc v$.
The preceding proposition identifies its unique geodesic with the full Hermitian Perron envelope. The base capacity endpoint limits give $L^1(Y)$ convergence. Since
\[
 \int_X|\pi^*(v-w)|\frac{\omega^n}{2n}
 =\int_Y|v-w|\kappa^d,
\]
the same holds on $X$; every smooth positive $\mathrm dV$ has bounded density ratio with $\omega^n$ on compact $X$. For smooth $v$, the constant translation is a base geodesic with speed $|c|$. Smooth metric approximation and translation by the same constant extend the distance formula to every $v\in\mathcal E^2$.
\end{proof}

For an absolutely continuous curve $\gamma:[0,1]\to\mathcal B_2$, define its metric action by
\[
 \mathcal A_B(\gamma)=\frac12\int_0^1|\dot\gamma|_{d_B}^{\,2}\,\mathrm dt.
\]
Here $|\dot\gamma|_{d_B}$ is the metric derivative. The action is defined even when the potentials are unbounded; it does not require a Riemannian tensor at a singular potential.

\Needspace{12\baselineskip}
\begin{proposition}\label{prop:hopf-basic-action}\index{Action!basic Hopf potentials}
The unique action minimizer between two points of $\mathcal B_2$ is their constant-speed Perron segment, and its action is half their squared distance. On every smooth path in $\mathcal B_\infty$, $\mathcal A_B$ equals the raw action $\mathscr S$. Every smooth basic Perron segment with strictly positive spatial forms is stationary for $\mathscr S$ under all smooth fixed-endpoint variations on $X$.
\end{proposition}
\begin{proof}
The length inequality and Cauchy--Schwarz give
\[
 d_B(\gamma_0,\gamma_1)^2
 \leq\left(\int_0^1|\dot\gamma|_{d_B}\,\mathrm dt\right)^2
 \leq2\mathcal A_B(\gamma).
\]
Equality holds for the constant-speed geodesic. Equality for a connecting curve forces constant speed and equality between length and endpoint distance, hence forces a geodesic; uniqueness follows from the CAT$(0)$ property.

For a smooth basic path $\pi^*v_s$, put $h=v_{s+\delta}-v_s$. On its short base Perron segment, parametrized by $0\leq\tau\leq1$, the chord bounds above, and
\[
 v_s+\tau h-C\delta^2\tau(1-\tau)
\]
bounds below. Indeed, the spatial forms along the chord have a uniform positive lower bound, $\mathrm dh=O(\delta)$, and \cref{prop:hermitian-schur-geodesic} makes the displayed path a subgeodesic for a fixed sufficiently large $C$. The right initial velocity of the Perron segment is therefore $h+O(\delta^2)$ uniformly. The classical smooth-endpoint K\"ahler distance formula gives
\[
 \lim_{\delta\to0}\frac{d_B(\pi^*v_s,\pi^*v_{s+\delta})^2}{\delta^2}
 =2n\int_Y\dot v_s^{\,2}(\kappa+\ddc v_s)^d.
\]
Integration proves $\mathcal A_B=\mathscr S$ along this path.

For stationarity, let $L(w)$ be the first variation at a smooth basic Perron path. Uniform spatial positivity allows every smooth endpoint-zero variation. Group invariance of the action and of the path gives $L(g^*w)=L(w)$ for $g\in E$. Average over Haar probability measure to obtain
\[
 L(w)=L\left(\int_Eg^*w\,\mathrm dg\right)=L(\pi^*q)
\]
for a smooth endpoint-zero base variation $q$. Fiber integration reduces this to $2n$ times the closed-base first variation. With $\eta_t=\kappa+\ddc v_t$, integration by parts gives that variation as
\[
 \int_0^1\int_Yq_t\left[-\ddot v_t\eta_t^d
 +d\,\mathrm d\dot v_t\wedge\dc\dot v_t\wedge\eta_t^{d-1}\right]\mathrm dt.
\]
The bracket vanishes by the smooth base homogeneous equation. Thus $L(w)=0$ for every variation on $X$.
\end{proof}

The metric action minimizes over curves in $\mathcal B_2$. The stationarity assertion allows all smooth variations. It does not imply minimization over all smooth Hermitian paths. Likewise, the initial-velocity formula used above has smooth endpoints; it is not a definition of the metric on all of $\mathcal B_2$.

\section{Negative directions of the raw action}
\label{sec:hopf-action-saddles}

On this Hopf background, constant Perron translations are stationary. Radial variations detect the failure of minimization. Every smooth two-periodic function of $r$ defines a smooth function on $X$, and for a radial path $u(t,r)$,
\[
 \omega+\ddc u=(1+u_r)\alpha_n+(1+u_{rr})\beta_n.
\]
For $\langle F\rangle=\frac12\int_0^2F(r)\,\mathrm dr$, fiber integration therefore gives
\begin{equation}\label{eq:hopf-radial-action}
 \mathscr S(u)=n\int_0^1
 \left\langle u_t^2(1+u_r)^d(1+u_{rr})\right\rangle\mathrm dt.
\end{equation}
Changing the local base section shifts the start of the full period and leaves the average unchanged.

\Needspace{12\baselineskip}
\begin{theorem}\label{thm:hopf-action-saddle}\index{Action!Hopf saddle}
For every nonzero constant $c$, the full Perron segment $u_t=ct$ is stationary for the raw action, whose Hessian has an infinite-dimensional negative subspace. The infimum of raw lengths over all smooth strictly positive paths from $0$ to $c$ is strictly smaller than $d_B(0,c)=\sqrt{2n}|c|$.
\end{theorem}
\begin{proof}
The constant path is basic, equals the chord, and is the full Perron segment by \cref{cor:translations-gauduchon}. Its stationarity also follows from \cref{prop:hopf-basic-action}.
Let $w(t,r)$ be smooth, two-periodic in $r$, and zero at $t=0,1$. The paths $ct+\varepsilon w$ have positive spatial forms for sufficiently small $|\varepsilon|$. Expanding \eqref{eq:hopf-radial-action} to second order gives
\[
 2n\int_0^1\left\langle
 w_t^2+2cw_t(dw_r+w_{rr})
 +c^2\left(\frac{d(d-1)}2w_r^2+dw_rw_{rr}\right)
 \right\rangle\mathrm dt.
\]
The integral of $w_rw_{rr}$ over a period vanishes, while
$\int_0^1\langle w_tw_{rr}\rangle\mathrm dt
=-\frac12[\langle w_r^2\rangle]_0^1=0$.
Completing the square gives the Hessian
\begin{equation}\label{eq:hopf-action-hessian}
 2n\int_0^1\left\langle
 (w_t+dcw_r)^2-\frac{d(d+1)}2c^2w_r^2
 \right\rangle\mathrm dt.
\end{equation}
For a positive integer $k$, take
\[
 w_k(t,r)=\sin(\pi t)\cos\bigl(\pi k(r-dct)\bigr).
\]
Then $w_{k,t}+dcw_{k,r}=\pi\cos(\pi t)\cos(\pi k(r-dct))$, and \eqref{eq:hopf-action-hessian} is
\[
 \frac{n\pi^2}2\left(1-\frac{d(d+1)}2c^2k^2\right).
\]
It is negative for all sufficiently large $k$. Distinct modes and their derivatives are orthogonal under the $r$ average, so the quadratic form is negative on every nonzero finite combination of these modes. Each combination gives a positive path for sufficiently small variation parameter and has action strictly below $nc^2$. Its raw length satisfies $L^2\leq2\mathscr S$ on the unit interval, so it is strictly below $\sqrt{2n}|c|$.
\end{proof}

These competing paths need not be subgeodesics. The strict length inequality does not establish definiteness or completeness of the unrestricted length infimum. For $n=3$ and $c=1$, a concrete path is
\[
 u_\varepsilon(t,r)=t+\varepsilon\sin(\pi t)\cos\bigl(\pi(r-2t)\bigr).
\]
Its spatial coefficients are at least $1-\pi|\varepsilon|$ and $1-\pi^2|\varepsilon|$. Its exact action is
\begin{equation}\label{eq:hopf-saddle-polynomial}
 \mathscr S(u_\varepsilon)
 =3-\frac{3\pi^2}2\varepsilon^2+\frac{111\pi^4}{64}\varepsilon^4.
\end{equation}
Indeed, put $f=\sin(\pi t)$, $p=f'$, and $\vartheta=\pi(r-2t)$. Substitute
$u_t=1+\varepsilon(p\cos\vartheta+2\pi f\sin\vartheta)$,
$u_r=-\varepsilon\pi f\sin\vartheta$, and $u_{rr}=-\varepsilon\pi^2f\cos\vartheta$ into \eqref{eq:hopf-radial-action}. Odd powers have zero periodic average. The averaged quadratic coefficient is $(p^2-3\pi^2f^2-2\pi^2fp)/2$, and the quartic coefficient is $\pi^2f^2(p^2+6\pi^2fp+12\pi^2f^2)/8$. The identities
\[
 \int_0^1fp\,\mathrm dt=\int_0^1f^3p\,\mathrm dt=0,
 \quad \int_0^1f^2\,\mathrm dt=\frac12,
 \quad \int_0^1p^2\,\mathrm dt=\frac{\pi^2}2,
 \quad \int_0^1f^4\,\mathrm dt=\frac38,
 \quad \int_0^1f^2p^2\,\mathrm dt=\frac{\pi^2}8
\]
give \eqref{eq:hopf-saddle-polynomial}. Taking $\varepsilon=1/100$ gives a positive path with action $2.9985212487\ldots<3$.

\section{Explicit basic geodesics and singular potentials}

Diagonal projective automorphisms give smooth nonconstant geodesics. For real $a_0,\ldots,a_{n-1}$, set
\[
 v_t([z])=\log\frac{\sum_{j=0}^{n-1}\mathrm e^{ta_j}|z_j|^2}
 {\sum_{j=0}^{n-1}|z_j|^2}.
\]
The family is spatially nonconstant when the $a_j$ are not all equal.

\Needspace{10\baselineskip}
\begin{proposition}\label{prop:hopf-diagonal-geodesic}\index{Geodesic!diagonal projective family}
The pullback $\pi^*v_t$ is the full Hermitian Perron segment and is stationary under all smooth fixed-endpoint variations. Its distance and action are
\[
 d_B(0,\pi^*v_1)^2=\frac{2[(\sum_ja_j)^2+\sum_ja_j^2]}{n+1},
 \qquad
 \mathscr S(\pi^*v)=\mathcal A_B(\pi^*v)
 =\frac{(\sum_ja_j)^2+\sum_ja_j^2}{n+1}.
\]
\end{proposition}
\begin{proof}
On the strip, the holomorphic map
$F([z],s)=[\mathrm e^{a_0s/2}z_0:\cdots:\mathrm e^{a_{n-1}s/2}z_{n-1}]$
satisfies $\operatorname{pr}_Y^*\kappa+\ddc V=F^*\kappa$. The imaginary-time phases are unitary, so $V$ depends only on $\Real s$. This form has rank $d$, strictly positive spatial slices, and zero $(d+1)$st power. The classical bounded K\"ahler Dirichlet uniqueness identifies it with the base Perron segment; \cref{prop:hopf-basic-descent,prop:hopf-basic-action} give the full segment and its stationarity.

For $x_j=|z_j|^2/\sum|z_j|^2$, the initial velocity is $\sum_ja_jx_j$. A standard complex Gaussian vector has independent exponential squared coordinates, and its projective direction has distribution $\kappa^d$. Dividing those coordinates by their sum gives the uniform simplex distribution, with moments
\[
 \int_Yx_i^2\kappa^d=\frac2{n(n+1)},
 \qquad \int_Yx_ix_j\kappa^d=\frac1{n(n+1)}\quad(i\ne j).
\]
The smooth distance formula and \eqref{eq:hopf-basic-metric} give the claimed squared distance; the action is half of it.
\end{proof}

For instance, $n=3$ and $(a_0,a_1,a_2)=(1,-1,0)$ give squared distance one and action $1/2$.

The completed space also contains potentials with different ordinary singularity types. Put $o=[1:0:\cdots:0]\in Y$ and
\[
 \psi=\log\frac{\sum_{j=1}^d|z_j|^2}{\sum_{j=0}^d|z_j|^2},
 \qquad v_a=1-(1-\psi)^a,\qquad0<a<1,
\]
with the canonical value $v_a(o)=-\infty$.

\Needspace{12\baselineskip}
\begin{proposition}\label{prop:hopf-singular-basic-energy}\index{Finite energy!radial exponent threshold}
Every $v_a$ has non-pluripolar mass one, and
\[
 v_a\in\mathcal E^2(Y,\kappa)
 \quad\Longleftrightarrow\quad a<\frac d{d+2}.
\]
At equality its squared energy diverges logarithmically. The pullback has mass $2n$ and pole set $\pi^{-1}(o)$. If $0<a<b<d/(d+2)$, the two pullbacks have different ordinary singularity types and finite positive $d_B$ distance.
\end{proposition}
\begin{proof}
In the chart $z_0=1$, write $s=\log|z'|^2$, $L=\log(1+\mathrm e^s)$, $P=\mathrm e^s/(1+\mathrm e^s)$, and $T=1+L-s\geq1$. The local potential $\Phi=L+1-T^a$ has derivative
\[
 A=P+a(1-P)T^{a-1},\qquad
 A'=P'(1-aT^{a-1})+a(1-a)(1-P)^2T^{a-2}>0.
\]
Both $A$ and $A'$ are positive, so the radial complex Hessian is positive. Its bounded-above local potential extends plurisubharmonically across $o$ by \cref{lem:qpsh-removable-analytic}. The function extends smoothly across the hyperplane at infinity. Its bounded plurisubharmonic truncations decrease to a nontrivial canonical potential: near $o$, $|v_a|=O((1-s)^a)$ is integrable against radial smooth volume $O(\mathrm e^{ds})\,\mathrm ds$.

The radial mass flux is $Q(s)=A(s)^d$. Its sphere coefficient is one, since $\dc\log|z'|^2$ integrates to one on the Hopf circle and $\ddc\log|z'|^2$ induces mass-one Fubini--Study on $\mathbb P^{d-1}$. Stokes on an annulus gives mass $Q(s_2)-Q(s_1)$. Since $Q(-\infty)=0$ and $Q(+\infty)=1$, finite-level locality and annular exhaustion give non-pluripolar mass one. The pole has no atom, and smoothness gives no mass on the hyperplane at infinity.

As $s\to-\infty$,
\[
 |v_a|^2\sim T^{2a},\qquad
 Q'(s)\sim d a^d(1-a)T^{-d(1-a)-1},\qquad T\sim1-s.
\]
The squared energy therefore converges exactly when $2a-d(1-a)<0$. At equality the integrand has order $T^{-1}$, giving logarithmic divergence. All other compact regions have smooth potentials and measures. Fiber integration gives the pullback assertions. Finally, $v_a-v_b=T^b-T^a\to+\infty$ at $o$, so mutual bounded comparison fails. Both points lie in the completed metric space and are distinct, giving finite positive distance.
\end{proof}

\section{Model projection and test curves}
\label{sec:hopf-basic-models}

The distance just constructed retains information about the potentials that their ordinary model projection forgets. This distinction connects the example with the singularity spaces studied earlier and with the Legendre construction in the next chapter.

\Needspace{12\baselineskip}
\begin{proposition}\label{prop:hopf-basic-model-projection}\index{Model envelope!basic Hopf potentials}
For every $v\in\PSH(Y,\kappa)$,
\[
 P_\omega[\pi^*v]=\pi^*P_\kappa[v].
\]
Consequently every point of $\mathcal B_2$ has ordinary model projection zero. Pullback commutes with the forward and backward Legendre transforms of base ordinary model test curves. If a base Perron ray has every finite slice, including its initial potential, in $\mathcal E^2(Y,\kappa)$, each finite segment of its pullback is a $d_B$ geodesic and minimizes the metric action with its prescribed endpoints.
\end{proposition}
\begin{proof}
For finite $C$, the obstacle $\pi^*\min(v+C,0)$ and the full competitor family are invariant under $E$. The same argument as in \cref{prop:hopf-basic-descent} gives
\[
 P_\omega(\pi^*\min(v+C,0))
 =\pi^*P_\kappa(\min(v+C,0)).
\]
Increasing limits and their upper semicontinuous regularizations commute with pullback in submersion charts, proving the model identity. The classical full-mass characterization on $Y$ gives $P_\kappa[v]=0$ for every $v\in\mathcal E^2$, hence the assertion for $\mathcal B_2$.

Pullback preserves parameterwise concavity, order, ordinary-model normalization and cutoff values of a base model test curve. It commutes with pointwise infima, and the submersion property makes it commute with upper semicontinuous regularization of suprema. Thus it commutes with both Legendre transforms. Finally, the Perron descent identifies every finite ray segment with the full Hermitian segment; its finite-energy endpoints and \cref{thm:hopf-basic-completion,prop:hopf-basic-action} give the geodesic and minimizing assertions.
\end{proof}

For ordinary-equivalent endpoints, \cref{thm:ordinary-type-perron} already gives their common type throughout the segment. The basic mass formula computes the mass of a chosen basic potential; it does not compute the lower or upper volume of its entire ordinary type, which also includes nonbasic potentials. A separating complete metric on all ordinary models, with compatibility under modifications and changes of background, remains open.
